\documentclass[11pt,letterpaper]{article}
\usepackage[utf8]{inputenc}
\usepackage{geometry}
\usepackage{setspace}

\geometry{letterpaper, margin=1in}
\doublespacing

\setlength{\parindent}{0.5in}
\setlength{\parskip}{0pt}

\begin{document}

\thispagestyle{empty}

\begin{center}
    \textbf{\large SPECIFICATION}
\end{center}

\vspace{0.4cm}

\begin{center}
    \textbf{\large OPTICALLY-ATTACHED STORAGE APPARATUS AND SYSTEM}
\end{center}

\vspace{0.4cm}

\section*{Cross-Reference to Related Applications}

This application claims the benefit of and priority to U.S. Patent Application No. 19/559,705, filed on March 6, 2026, the contents of which are hereby incorporated by reference in their entirety.

\section*{Background of the Invention}

The performance of mass storage devices, such as solid-state drives (SSDs) and hard disk drives (HDDs), has long been constrained by the data interfaces used to connect them to a computer system. Interfaces like SATA, SAS, and even the more modern NVMe are fundamentally serial or limited-parallel electrical standards. While the speed of the storage media themselves (e.g., NAND flash) has increased dramatically, the interface acts as a bottleneck, preventing the system from realizing the full read/write potential of the drive. This I/O bottleneck is a major limiting factor in high-performance computing, large-scale data analysis, and any application requiring the rapid movement of massive datasets.

\section*{Summary of the Invention}

The present invention is an apparatus for a mass storage device, referred to as Optically-Attached Storage (OAS), that eliminates the traditional interface bottleneck by using a massively parallel optical port for all data I/O. The invention is a direct application of the technology described in the parent application, No. 19/559,705, to the field of data storage.

The OAS device comprises a non-volatile storage medium, a storage controller, and a massively parallel optical interface composed of a two-dimensional array of photon emitters and a corresponding array of photon sensors. This optical port replaces conventional physical connectors. Data is read from or written to the storage medium by transmitting or receiving a sequence of two-dimensional frames of light, where each frame contains millions of parallel data points encoded in the light's characteristics.

This architecture allows the device to transfer data at rates of hundreds of gigabits or even terabits per second, a rate that is commensurate with the internal speed of the storage medium itself. The OAS drive is a foundational hardware component for building ultra-high-performance servers and is the key enabling technology for a new ecosystem of high-speed physical data logistics.

\section*{Statement of License Grant to the United States Government}

The inventor, Reinaldo Javier Aguiar Marcano, hereby grants to the United States Federal Government and its military and defense agencies a perpetual, royalty-free, non-exclusive license to practice and use the invention described herein for any and all governmental, defense, and non-commercial purposes. This license explicitly excludes any and all state and local governments.

\section*{Brief Description of the Drawings}

FIG. 1 is a diagram contrasting a conventional storage drive with its serial electrical interface against the Optically-Attached Storage (OAS) drive with its massively parallel optical interface.

FIG. 2 is a block diagram illustrating the internal components of the OAS drive.

\section*{Detailed Description of the Invention}

FIG. 1 provides a conceptual contrast between prior art and the present invention. A conventional storage drive (100) contains a storage medium (101) and a controller (102), but its data I/O is funneled through a narrow, serial electrical interface (103) like SATA or NVMe, creating a bottleneck. The Optically-Attached Storage (OAS) drive (110) of the present invention also contains a storage medium (111) and a controller (112). However, its data I/O is handled by a wide, massively parallel optical interface (113), which allows for unimpeded data flow at a much higher bandwidth.

Referring to FIG. 2, the internal architecture of the OAS drive (110) is shown. The core components are the non-volatile storage medium (111) (e.g., NAND flash chips), the storage controller (112), and the optical interface (113). The optical interface itself comprises a data-receiving array of photon sensors (201) and a data-transmitting array of photon emitters (202), allowing for bidirectional communication. When a `write' command is received from a host system via the sensor array (201), the controller (112) directs the data to be written to the storage medium (111). When a `read' command is received, the controller (112) retrieves the data from the medium (111) and sends it to the emitter array (202) to be transmitted to the host.

\end{document}
